\documentclass[11pt]{article}
\usepackage[margin=1in]{geometry}
\usepackage[T1]{fontenc}
\usepackage[utf8]{inputenc}
\usepackage{lmodern}
\usepackage{amsmath,amssymb,amsthm}
\usepackage{booktabs}
\usepackage{graphicx}
\usepackage{hyperref}
\usepackage{xcolor}
\usepackage{enumitem}
\usepackage{microtype}

\hypersetup{
  colorlinks=true,
  linkcolor=blue,
  citecolor=blue,
  urlcolor=blue
}

\title{When Forecast Accuracy Is Not Enough:\\Calibration-Aware Autoscaling for Bursty Serverless Workloads}
\author{Dang Le}
\date{\today}

\begin{document}
\maketitle

\begin{abstract}
Autoscaling policies in cloud systems are commonly either reactive or driven by point forecasts of future demand. Under bursty workloads, both approaches can be brittle: reactive methods respond too late, while point-forecast policies ignore predictive uncertainty and may underprovision during sudden demand spikes. We propose a calibration-aware autoscaling framework that converts predictive uncertainty into scaling decisions through an explicit cost-reliability objective. Rather than scaling to a fixed quantile, our method uses calibrated predictive intervals to balance provisioning cost, underprovisioning risk, and scaling instability. We evaluate the approach on the public Azure Functions Invocation Trace 2021, aggregated at the application level, against reactive, moving-average, point-forecast, and fixed-quantile baselines. In the current implementation, predictive methods reduce aggregate SLA violations substantially relative to reactive scaling, but burst-only windows remain difficult for all methods, which clarifies where calibration-aware decisions must improve further. Our central claim is that forecast accuracy alone is insufficient for robust autoscaling: calibration quality materially affects downstream scaling decisions.
\end{abstract}

\section{Introduction}
Cloud autoscaling is fundamentally a decision-making problem under uncertainty. In practice, many autoscaling policies are either reactive, triggered by observed threshold violations, or predictive, relying on point forecasts of future workload demand. Both paradigms struggle when workloads are bursty, non-stationary, or exhibit abrupt shifts. Reactive methods often respond too late, while point forecasts conceal uncertainty and encourage brittle scaling decisions.

This paper argues that forecast accuracy alone is not enough for robust autoscaling. What matters is whether predictive uncertainty is sufficiently calibrated to support downstream capacity decisions. A forecast with low point error but poorly calibrated uncertainty can still produce costly or unsafe provisioning outcomes. We therefore study calibration-aware autoscaling: a framework that explicitly integrates predictive uncertainty into the scaling policy through a cost-reliability objective.

\paragraph{Problem statement.}
Given historical application-level invocation traces and a short prediction horizon, can calibrated predictive uncertainty improve autoscaling decisions relative to reactive and point-forecast baselines under bursty serverless workloads?

\paragraph{Thesis.}
Calibration-aware autoscaling can reduce SLA violations and improve the cost-reliability trade-off under bursty serverless workloads, and the quality of uncertainty calibration is itself a key determinant of downstream scaling quality.

\section{Research Questions and Hypotheses}
\subsection{Research Questions}
\begin{enumerate}[label=RQ\arabic*., leftmargin=1.2cm]
    \item Can calibrated predictive uncertainty improve autoscaling decisions under bursty serverless workloads compared with reactive and point-forecast baselines?
    \item Does calibration quality materially affect downstream scaling quality, beyond standard forecasting accuracy metrics such as MAE or RMSE?
    \item Under which workload regimes (stable, periodic, bursty, drifting) does calibration-aware autoscaling provide the largest benefit?
\end{enumerate}

\subsection{Hypotheses}
\begin{enumerate}[label=H\arabic*., leftmargin=1.2cm]
    \item Calibration-aware autoscaling reduces SLA violations relative to reactive autoscaling under bursty serverless workloads.
    \item Calibration-aware autoscaling achieves a better cost-reliability trade-off than point-forecast-based autoscaling.
    \item Forecast accuracy alone is insufficient; poorly calibrated uncertainty can lead to inferior scaling decisions even when point forecast error is low.
\end{enumerate}

\section{Problem Formulation and Theoretical Motivation}
Let $H_t$ denote the observed workload history up to time $t$, and let $D_{t+h}$ denote the random required capacity at forecast horizon $h$. The autoscaling decision is to choose an integer capacity $c_t \in \mathbb{Z}_{\ge 0}$ before $D_{t+h}$ is observed.

The paper treats autoscaling as a decision problem under predictive uncertainty. A forecasting module produces either a predictive distribution or a point forecast with an uncertainty estimate for $D_{t+h}$. The policy then solves
\[
\mathcal{L}_t(c_t) =
\lambda_c \, c_t
\;+\;
\lambda_u \, \mathbb{E}\!\left[(D_{t+h} - c_t)_+\right]
\;+\;
\lambda_s \, |c_t - c_{t-1}|,
\]
where $(x)_+ = \max(x,0)$, $\lambda_c > 0$ is the unit provisioning cost, $\lambda_u > 0$ is the underprovision penalty weight, and $\lambda_s > 0$ penalizes oscillatory scaling.

This objective is intentionally simple. It is rich enough to connect predictive uncertainty to downstream decisions, while remaining interpretable and easy to evaluate on a public trace.

\paragraph{Proposition 1 (Quantile rule in the static case).}
Consider the special case $\lambda_s = 0$ and assume the demand distribution for $D_{t+h}$ is known exactly. Then any optimizer of
\[
\lambda_c \, c + \lambda_u \, \mathbb{E}\!\left[(D_{t+h} - c)_+\right]
\]
corresponds to a quantile decision rule. In particular, for the continuous idealization, the optimal capacity is an $\alpha$-quantile of $D_{t+h}$ with
\[
\alpha = 1 - \frac{\lambda_c}{\lambda_u},
\qquad \text{provided } 0 < \lambda_c < \lambda_u.
\]

\paragraph{Proof sketch.}
The derivative of the objective with respect to $c$ is
\[
\frac{d}{dc}\mathcal{L}(c)
=
\lambda_c - \lambda_u \, \mathbb{P}(D_{t+h} > c).
\]
At optimality, the tail probability balances marginal cost and marginal underprovision penalty, so the solution is a quantile of the demand distribution. This shows that fixed-quantile scaling is not arbitrary: it is the optimizer of a simplified asymmetric-cost decision problem.

\paragraph{Why this matters.}
Proposition 1 gives a clean theoretical interpretation of the fixed-quantile baseline and also clarifies the proposed method: calibration-aware scaling is a generalization of this quantile logic to a setting with uncertainty estimation errors and a temporal smoothness penalty.

\paragraph{Proposition 2 (Why calibration matters for decisions).}
Suppose a policy chooses capacity from a predictive quantile or predictive interval under the intention of enforcing a target tail risk level. If the predictive distribution is miscalibrated, then the realized exceedance probability
\[
\mathbb{P}(D_{t+h} > c_t)
\]
can differ substantially from the intended design level. Consequently, the actual cost-reliability operating point can drift away from the one implied by the policy weights or chosen quantile.

\paragraph{Interpretation.}
This proposition is conceptually simple but central to the paper. A fixed P95 policy is only meaningful if the estimated P95 behaves like a true 95th percentile under the data distribution. If calibration is poor, then even a seemingly conservative quantile rule can underperform in downstream scaling. This is the main reason the paper emphasizes calibration quality rather than forecast error alone.

\paragraph{Observation 1 (Role of the oscillation penalty).}
When $\lambda_s > 0$, the policy no longer reduces to a static quantile rule. The optimal decision becomes history dependent through $c_{t-1}$, which creates a direct trade-off between burst responsiveness and scaling stability. This is precisely where a calibration-aware decision policy differs from a naive fixed-quantile baseline.

\section{Method}
\subsection{Overview}
The proposed system has three layers:
\begin{enumerate}[leftmargin=1.2cm]
    \item \textbf{Forecasting layer:} predicts short-horizon demand and produces predictive uncertainty.
    \item \textbf{Calibration layer:} adjusts uncertainty estimates so that predictive intervals better match empirical coverage.
    \item \textbf{Decision layer:} converts the calibrated demand distribution into a scaling decision using an explicit cost-risk objective.
\end{enumerate}

\subsection{Forecasting Layer}
We will compare a small set of forecasters to avoid unnecessary complexity:
\begin{itemize}[leftmargin=1.2cm]
    \item Naive / moving-average baseline
    \item One point-forecast baseline based on linear autoregression
    \item One neural baseline (e.g., LSTM or TCN), only after simpler baselines are stable
    \item One probabilistic forecaster producing quantiles or predictive intervals
\end{itemize}

\subsection{Calibration Layer}
The first implementation uses a simple split-calibration procedure on a validation split. A point-forecast model is trained on the training partition, then its residuals on the validation partition are collected to form an empirical error distribution. These residuals are used in two ways:
\begin{itemize}[leftmargin=1.2cm]
    \item to construct calibrated predictive intervals via residual quantiles,
    \item to form a residual-based demand sample set for the downstream decision objective.
\end{itemize}
This gives a transparent calibration-aware method before introducing heavier probabilistic models.

\subsection{Decision Layer}
Let $D_t$ denote the predicted demand distribution at time $t$, and let $c_t$ denote the chosen capacity. The scaling decision minimizes a risk-sensitive objective:
\[
\mathcal{L}(c_t) = \lambda_1 \cdot \mathrm{Cost}(c_t)
+ \lambda_2 \cdot \mathbb{E}[\mathrm{UnderProvisionPenalty}(D_t, c_t)]
+ \lambda_3 \cdot |c_t - c_{t-1}|.
\]

This formulation balances:
\begin{itemize}[leftmargin=1.2cm]
    \item provisioning cost,
    \item penalty for underprovisioning,
    \item scaling instability or oscillation.
\end{itemize}

\subsection{Data-to-Decision Setup}
The empirical study uses the Azure Functions Invocation Trace 2021 as a concrete serverless case study. For each application, invocation start times are reconstructed from end timestamps and durations, then aggregated into 1-minute buckets. The phase-1 target is minute-level required capacity, defined from average concurrent demand within each minute. This choice keeps the decision variable close to a practical autoscaling proxy while remaining faithful to the public trace.

\subsection{Why this method is publishable}
The novelty is not merely ``use a better model.'' The paper focuses on three deeper claims:
\begin{enumerate}[leftmargin=1.2cm]
    \item uncertainty should be calibrated before it is trusted for scaling,
    \item decision quality depends on calibration, not only point accuracy,
    \item bursty regimes expose the weaknesses of reactive and point-forecast policies more clearly than aggregate metrics do.
\end{enumerate}

\section{Experimental Plan}
\subsection{Dataset Strategy}
Start with \textbf{one public workload trace} only. This keeps the scope realistic for a first paper. We use the \textbf{Azure Functions Invocation Trace 2021} as the sole dataset. The study aggregates invocations at the application level, uses 1-minute buckets, and predicts required capacity at horizons of 5, 10, and 15 minutes.

\subsection{Baselines}
We will compare against four baselines:
\begin{enumerate}[leftmargin=1.2cm]
    \item Reactive threshold autoscaling
    \item Moving-average forecast + scaling decision
    \item Linear autoregressive point-forecast scaling
    \item Fixed-quantile scaling (e.g., rolling P90 or P95)
\end{enumerate}

\subsection{Evaluation Metrics}
\paragraph{Forecasting metrics}
\begin{itemize}[leftmargin=1.2cm]
    \item MAE / RMSE
    \item Pinball loss (for quantile forecasts)
\end{itemize}

\paragraph{Calibration metrics}
\begin{itemize}[leftmargin=1.2cm]
    \item Empirical coverage
    \item Calibration gap
    \item Average interval width
\end{itemize}

\paragraph{Downstream scaling metrics}
\begin{itemize}[leftmargin=1.2cm]
    \item Provisioning cost
    \item SLA violation rate
    \item Underprovision ratio
    \item Scaling frequency / oscillation
\end{itemize}

\subsection{Core Experiments}
\begin{enumerate}[leftmargin=1.2cm]
    \item Full comparison across all methods on the test split.
    \item Burst-only evaluation on windows with sharp demand spikes.
    \item Ablation over forecast horizon, confidence level, and penalty weights.
    \item Calibration stress test: compare residual-calibrated intervals versus fixed-quantile decisions and miscalibrated uncertainty.
\end{enumerate}

\section{Results}
\subsection{Experimental Setup}
The current codebase has been executed end-to-end on the Azure Functions Invocation Trace 2021. The experimental setup uses:
\begin{itemize}[leftmargin=1.2cm]
    \item the top 10 active applications by invocation volume,
    \item 1-minute aggregation,
    \item a required-capacity proxy derived from average concurrent demand per minute as the main configuration,
    \item chronological train/validation/test splits with proportions 70/15/15,
    \item a 5-minute forecast horizon for the first round of experiments.
\end{itemize}

This implementation is intentionally simple and reproducible. It should be viewed as the first end-to-end empirical milestone rather than the final version of the method.

\subsection{Aggregate Comparison}
Table~\ref{tab:aggregate-results} summarizes the first aggregate comparison on the test split.

\begin{table}[h]
\centering
\caption{Preliminary aggregate comparison on the Azure Functions 2021 test split. Lower is better for cost, SLA violation, and oscillation.}
\label{tab:aggregate-results}
\begin{tabular}{lccc}
\toprule
Method & Mean Cost & SLA Violation Rate & Mean Oscillation \\
\midrule
Reactive & 0.626 & 0.0318 & 0.1566 \\
Moving average & 0.543 & 0.0197 & 0.0225 \\
Point forecast & 1.044 & 0.0114 & 0.0754 \\
Fixed quantile (0.95) & 1.067 & 0.0205 & 0.0480 \\
Risk-calibrated & 1.038 & 0.0115 & 0.0678 \\
\bottomrule
\end{tabular}
\end{table}

Two observations already stand out. First, reactive scaling is clearly inferior on aggregate reliability: its SLA violation rate is roughly three times that of the point-forecast and calibration-aware methods. Second, the risk-calibrated policy is competitive with the point-forecast baseline on aggregate performance, yielding slightly lower mean cost and lower oscillation, but not yet a decisive SLA improvement.

\subsection{Calibration Metrics}
For the risk-calibrated implementation with nominal coverage $0.90$, the empirical coverage on the test split is $0.8793$, with calibration gap $0.0207$ and average interval width $0.8972$. This is close enough to support the decision layer meaningfully, but it also shows that calibration is not yet perfect.

\subsection{Burst-Only Evaluation}
Aggregate performance is not the full story. We therefore isolate burst windows using a rolling-baseline detector and compare methods only on those windows. In the current test split, burst windows make up only about $1.23\%$ of all decision steps, but they are much harder than the average case.

\begin{table}[h]
\centering
\caption{Preliminary burst-only comparison. Lower is better for cost, SLA violation, and oscillation.}
\label{tab:burst-results}
\begin{tabular}{lccc}
\toprule
Method & Mean Cost & SLA Violation Rate & Mean Oscillation \\
\midrule
Reactive & 4.881 & 0.7665 & 3.4229 \\
Moving average & 0.546 & 0.9604 & 0.0617 \\
Point forecast & 1.013 & 0.7885 & 0.0396 \\
Fixed quantile (0.95) & 0.925 & 0.9207 & 0.0044 \\
Risk-calibrated & 1.009 & 0.7885 & 0.0352 \\
\bottomrule
\end{tabular}
\end{table}

These results sharpen the paper's motivation. All current methods still fail badly on burst windows, including the calibration-aware one. However, they fail in different ways: reactive scaling becomes extremely expensive and unstable, while smoothing and fixed-quantile policies remain cheap but miss bursts almost completely. The point-forecast and risk-calibrated policies are the strongest among the current methods on burst windows, yet their SLA violation rates remain far too high. This suggests that the main research opportunity lies not in squeezing small gains from average-case metrics, but in improving burst sensitivity without giving up reproducibility.

\subsection{Discussion of the First Round}
The first empirical round supports several parts of the project's thesis:
\begin{itemize}[leftmargin=1.2cm]
    \item reactive autoscaling is not competitive on this trace,
    \item fixed-quantile scaling is not automatically safer than a calibrated decision rule,
    \item burst-only analysis reveals weaknesses that aggregate metrics can hide.
\end{itemize}

At the same time, the current evidence does \emph{not} yet justify a strong claim that the calibration-aware method clearly dominates point-forecast scaling. The next iteration should therefore focus on improving burst-sensitive decision behavior rather than adding a more complicated forecasting model too early.

\subsection{Capacity Proxy Ablation: Average vs Peak}
We also tested a sharper capacity proxy based on per-minute peak concurrency. This is a natural alternative because it should be more sensitive to bursts than average concurrency. However, the first ablation shows that the peak proxy makes the task substantially harsher for all methods and does not currently improve the downstream behavior of the calibration-aware policy.

\begin{table}[h]
\centering
\caption{Risk-calibrated policy under two capacity proxies. The average-concurrency proxy remains the main setting because it gives a better cost-reliability trade-off in the current implementation.}
\label{tab:proxy-ablation}
\begin{tabular}{lcccccc}
\toprule
Proxy & Mean Cost & SLA Viol. & Oscillation & Coverage & Burst Fraction & Burst SLA \\
\midrule
Average concurrency & 1.038 & 0.0115 & 0.0678 & 0.8793 & 0.0123 & 0.7885 \\
Peak concurrency & 1.179 & 0.0301 & 0.2195 & 0.9078 & 0.0225 & 0.8558 \\
\bottomrule
\end{tabular}
\end{table}

This ablation changes the interpretation of the next steps. The peak proxy does improve empirical coverage slightly, but it worsens cost, aggregate SLA violations, and burst-only SLA violations for the current policy. In other words, simply making the target sharper is not enough. For this paper, the average-concurrency proxy should remain the primary setup, while the peak proxy serves as a useful negative-result ablation that highlights the need for better burst-sensitive decision logic.

\section{Expected Contributions}
\begin{enumerate}[leftmargin=1.2cm]
    \item A calibration-aware autoscaling framework for bursty workloads.
    \item Evidence that calibration quality directly affects autoscaling quality.
    \item A reproducible empirical comparison against standard autoscaling baselines.
    \item A public research artifact: paper, code, and plots.
\end{enumerate}

\section{Concrete Implementation Roadmap}
\subsection{Phase 1: Scope Lock (Days 1--3)}
\begin{itemize}[leftmargin=1.2cm]
    \item Finalize title, problem statement, and success criteria.
    \item Lock Azure Functions 2021 as the only dataset.
    \item Define the exact prediction target as minute-level required capacity and set horizons to 5, 10, and 15 minutes.
    \item Create the repository structure: \texttt{data/}, \texttt{src/}, \texttt{experiments/}, \texttt{plots/}, \texttt{paper/}.
\end{itemize}

\subsection{Phase 2: Data Pipeline (Days 4--7)}
\begin{itemize}[leftmargin=1.2cm]
    \item Load and inspect the trace.
    \item Reconstruct invocation start timestamps and aggregate by application.
    \item Create train/validation/test splits by time.
    \item Build a sliding-window time-series dataset.
    \item Define a burst-detection rule for analysis.
\end{itemize}

\subsection{Phase 3: Baselines (Week 2)}
\begin{itemize}[leftmargin=1.2cm]
    \item Implement reactive threshold scaling.
    \item Implement smoothing-based forecast baseline.
    \item Implement point-forecast baseline.
    \item Run the first end-to-end experiment and verify that metrics and plots work.
\end{itemize}

\subsection{Phase 4: Proposed Method (Weeks 3--4)}
\begin{itemize}[leftmargin=1.2cm]
    \item Implement split-residual calibration for predictive intervals.
    \item Implement the risk-calibrated decision policy.
    \item Compare against fixed-quantile scaling.
    \item Add a heavier probabilistic forecaster only if the simpler pipeline is stable.
\end{itemize}

\subsection{Phase 5: Evaluation and Writing (Weeks 5--8)}
\begin{itemize}[leftmargin=1.2cm]
    \item Run core experiments.
    \item Produce the main figures and result tables.
    \item Write Introduction, Method, Experimental Setup, and Results.
    \item Release code, finalize the paper PDF, and prepare an arXiv or workshop-ready version.
\end{itemize}

\section{Immediate Next Actions}
\begin{enumerate}[leftmargin=1.2cm]
    \item Keep the average-concurrency proxy as the main experimental setting and report peak concurrency only as an ablation.
    \item Explore an intermediate proxy, such as a within-minute high quantile, only if it improves burst behavior without collapsing the aggregate trade-off.
    \item Improve the decision layer for burst windows before adding a heavier forecasting model.
    \item Add one more figure-level case study that contrasts a burst episode under reactive, point-forecast, and risk-calibrated decisions.
\end{enumerate}

\section{Limitations and Risk Control}
This paper intentionally avoids reinforcement learning, large-scale cluster simulation, and broad cloud optimization claims. The goal is a focused, reproducible first publication rather than a large, multi-component systems project. The main risk is scope creep; the primary mitigation is to commit to one dataset, a small set of baselines, and a narrow decision objective.

\section{Conclusion}
This project treats autoscaling as a risk-sensitive decision problem rather than a pure forecasting problem. Its central claim is simple but meaningful: uncertainty is useful for scaling only when it is calibrated well enough to support downstream decisions. If validated empirically, this insight would make the project strong enough for a first public paper while remaining realistic for a focused student effort.

\end{document}
